\section{Sharp grid transfer and certified coarsening}\label{sec:transfer}

\subsection{Supported rounding preserves the switch budget}

Integral prefix rounding and support-preserving matching constructions are
established tools. In particular, \citet[Corollary~2.7]{bestehorn2021vanishing}
prove a uniform-mesh error bound with positive-support restrictions. We give
the short flow construction needed here, then use temporal support to
preserve the switch budget of a supplied integer schedule.

\begin{theorem}[Uniform-grid transfer]\label{thm:sharp-transfer}
Let $W$ be a cumulative integer schedule with finitely many switches on
$[0,T]$, and let $\mathcal T=(0,h,\ldots,Nh)$ with $T=Nh$ and $h>0$.
There is a grid schedule $V$ with no more switches than $W$ such that
\begin{equation}\label{eq:schedule-transfer}
 \norm{V-W}_\infty<h.
\end{equation}
Its word is obtained from a chronological subsequence of the original block
word by merging adjacent equal labels. With rational original switch times
and $h$, it can be constructed in polynomial bit complexity in the explicit
grid length, number of modes, and input size.
\end{theorem}
\begin{proof}
Set $a_{ji}=h^{-1}(W_i(jh)-W_i((j-1)h))$ and
$S_{ji}=\sum_{q=1}^j a_{qi}$. Each row of $a$ is a simplex vector. Create
a supply-one node for each cell and a chronological chain for each mode.
The cell node may send flow to mode $i$'s chain at that cell only when
$a_{ji}>0$. The edge carrying mode $i$'s accumulated flow after cell $j$
has integer lower and upper capacities $\lfloor S_{ji}\rfloor$ and
$\lceil S_{ji}\rceil$; final chain edges enter a sink of demand $N$.
The flows $a_{ji}$ and $S_{ji}$ show fractional feasibility. Integral
capacitated-flow feasibility therefore gives assignments $y_{ji}\in\{0,1\}$
with exactly one per cell, supported where $a_{ji}>0$, and
\[
 \lfloor S_{ji}\rfloor\le\sum_{q=1}^j y_{qi}
 \le\lceil S_{ji}\rceil.
\]
Their distance from $S_{ji}$ is strictly below one, including zero distance
when $S_{ji}$ is integral. Hence the endpoint errors are strictly below $h$.
Within a cell, the selected component of $V-W$ is nondecreasing and every
other component is nonincreasing. The finitely many endpoint bounds prove
\eqref{eq:schedule-transfer} on the whole interval.

For each cell choose an interior time, away from the original switches,
at which the original schedule occupies the selected mode. Positive cell
occupation guarantees such a time. These times are strictly increasing, so
their original block indices are nondecreasing. Remove repeated selections
of the same original block. The selected word is then a subsequence of the
original word; merging adjacent equal labels only decreases its length.
It has no more changes than the original word. The network has polynomial
size and integer capacities at most $N$; standard integral flow algorithms
and exact rational occupation calculations give the computational claim.
\end{proof}

\begin{corollary}[Instance and minimax transfer]\label{cor:opt-transfer}
For every $A\in\Acal_n(T)$, switch budget $s\ge0$, and uniform grid of
width $h$,
\begin{equation}\label{eq:instance-transfer}
 \OPT_s(A)\le\OPT_s^{\mathcal T}(A)<\OPT_s(A)+h.
\end{equation}
The minimax values, with the input classes in~\eqref{eq:minimax-values}
and~\eqref{eq:grid-definitions}, satisfy
\begin{equation}\label{eq:minimax-transfer}
 F_{n,s}(T)\le F_{n,s}^{\mathcal T}<F_{n,s}(T)+h.
\end{equation}
\end{corollary}
\begin{proof}
Restriction of the schedule set gives the first instance inequality. Apply
Theorem~\ref{thm:sharp-transfer} to an attained continuous optimum and use
the triangle inequality for the strict upper bound. The first minimax
inequality is~\eqref{eq:minimax-grid-domination}, which uses exact cell
averaging. For the other inequality, take a grid minimax maximizer $A_*$,
whose existence follows from Proposition~\ref{prop:compactness}. Then
$F_{n,s}^{\mathcal T}=\OPT_s^{\mathcal T}(A_*)<\OPT_s(A_*)+h
\le F_{n,s}(T)+h$. Attainment is essential to this argument for strictness;
a supremum of arbitrary pointwise strict inequalities need not remain strict.
\end{proof}

The construction transfers a given schedule; it does not find a continuous
optimum. It preserves neither arbitrary minimum dwell times nor arbitrary
transition restrictions, since skipping a mode can create a forbidden
transition. Theorem~\ref{thm:sharp-transfer} needs a uniform grid: unequal cell
lengths do not provide the unit assignment capacities used in the proof.

\begin{proposition}[Binary transfer on a nonuniform grid]\label{prop:binary-transfer}
For two modes on any grid, a finitely switching schedule has a grid
replacement with no more switches and discrepancy at most
$\bar\Delta/2$. This constant is sharp for transfer of a supplied schedule.
\end{proposition}
\begin{proof}
Let $D=\bar\Delta$ and maintain the mode-one endpoint discrepancy
$e=V_1-W_1$ in $[-D/2,D/2]$. On a cell of length $d\le D$, let $a$ be
the original mode-one occupation. If $a=0$ select mode two; if $a=d$
select mode one. These supported choices leave $e$ unchanged. Otherwise,
both modes occur, and the two next discrepancies are $e-a$ and $e+d-a$.
The first is at most $D/2$, the second at least $-D/2$, and their distance
is $d\le D$. They cannot both lie outside the permitted interval, so choose
one inside. The other component's discrepancy is its negative. Cellwise
monotonicity and the support chronology in
Theorem~\ref{thm:sharp-transfer} finish the proof. On a single cell of length
$D$, an original schedule occupying each mode for half the cell has
discrepancy $D/2$ from either possible constant grid schedule.
\end{proof}

\begin{proposition}[Sharp universal coefficient]\label{prop:transfer-sharpness}
The coefficient one in~\eqref{eq:instance-transfer} cannot be reduced
uniformly over mode counts and budgets, even for grid-constant relaxed inputs
and $1\le s\le N-2$.
\end{proposition}
\begin{proof}
Use unit cells, $n\ge2$, $N=n+1$, and $s=n-1=N-2$. Let the relaxed
control be uniform on $[0,1]$ and equal to mode $n$ afterward. Every grid
schedule incurs error $1-1/n$ at time one in its first selected mode. The
constant mode-$n$ schedule attains that error, so its grid optimum is
$1-1/n$. A continuous schedule visits modes $1,\ldots,n$, each for time
$1/n$ in the first cell, then remains in mode $n$. It has $n-1$ switches.
All errors vanish at time one; within the first cell the largest absolute
error is $(n-1)/n^2$. Thus
\[
 \OPT_s^{\mathcal T}(A)-\OPT_s(A)
 \ge(1-1/n)-(n-1)/n^2=(1-1/n)^2\longrightarrow1.
\]
Scaling proves the assertion for any mesh width $h$.
\end{proof}

Repeating the within-cell cycle $m$ times, taking $N=nm+1$ and
$s=nm-1$, gives the additional family with continuous competitor error
$(n-1)/(n^2m)$ and gap at least $(1-1/n)(1-1/(nm))$.
Already $n=4,m=1$ gives a gap of at least $9/16>1/2$.
This is an instance-wise counterexample to the half-mesh justification
immediately preceding Conjecture~1 in
\citet[p.~615]{sager-zeile2021}, also present in
\citet[p.~139]{zeile2021thesis}. Those passages motivate a conjecture and
are not stated as proved transfer theorems. The example does not determine
the difference between two minimax values. For one switch,
Corollary~\ref{cor:one-switch-grid} gives the stronger nearest-boundary
half-mesh estimate.

\subsection{An additive certificate from exact coarse optimization}

\begin{theorem}[Certified coarsening]\label{thm:certified-coarsening}
Let $A\in\Acal_n(T)$ and fix a uniform coarse grid of $M\ge1$ cells,
$h=T/M$. Compute its exact grid optimum $U_h$ and an attaining schedule
$V_h$ using the exact cumulative values $A_i(jh)$. Then
\begin{equation}\label{eq:coarsening-certificate}
 U_h-h<\OPT_s(A)\le U_h=D(A,V_h).
\end{equation}
In particular, for $0<\varepsilon\le1$ and $M=\lceil1/\varepsilon\rceil$,
this returns a feasible schedule whose suboptimality is strictly below
$\varepsilon T$. If $A$ is specified by rational piecewise-constant rates
on an $N$-cell input grid, the coarse data can be formed in
$O(n(N+M))$ rational operations, even when the grids are not aligned.
For each fixed $s$, the resulting additive approximation algorithm has
polynomial bit complexity in the input length and $1/\varepsilon$.
\end{theorem}
\begin{proof}
Endpoint monotonicity evaluates $D(A,V_h)$ exactly from the coarse cumulative
values; the rates need not be constant in coarse cells.
Equation~\eqref{eq:instance-transfer} proves the interval. Since
$h\le\varepsilon T$, it also proves the additive error claim.

For piecewise-constant input, traverse the sorted input and coarse endpoints
together. On each segment of their common refinement, add its length times
the known input rate to the current coarse mass. The common refinement has
at most $N+M-1$ cells, so this takes $O(n(N+M))$ operations. This integration
preserves the original $A_i(jh)$; merely grouping whole fine cells would be
incorrect for nonaligned grids. Apply Theorem~\ref{thm:fixed-budget} with
$k=\min(s+1,M)$, or the faster one-switch specialization when applicable.
The arithmetic bound is
\begin{equation}\label{eq:coarsening-complexity}
 O\left(n(N+M)+\binom{M-1}{k-1}n(k2^k+3^k)\right).
\end{equation}
All integrations and subsequent computations use rationals of polynomial
bit length. For fixed $s$, the bound is polynomial in $N,n,M$ and input
bit length. More generally its dependence on $(s,\lceil1/\varepsilon\rceil)$
can be placed in a parameter-only factor multiplying a polynomial in the
input length.
\end{proof}

Some special cases admit sharper bounds. For $s=1$ on any grid $\mathcal T$,
move an attained continuous optimum's single switch to the nearer endpoint
of its containing cell. The time displacement, and hence the change in
cumulative occupation, is at most half that cell's length. Thus, with
$U_{\mathcal T}=\OPT_1^{\mathcal T}(A)$,
\begin{equation}\label{eq:one-switch-arbitrary-grid-certificate}
 U_{\mathcal T}-\bar\Delta/2\le\OPT_1(A)\le U_{\mathcal T}.
\end{equation}
This holds for every measurable input and every mode count in the model;
it does not assert attainment of the gap $\bar\Delta/2$ between optima.
For a uniform grid the estimate is $h/2$. For $n=2$ and any $s$,
Proposition~\ref{prop:binary-transfer} gives the same estimate. In either case,
\[
 U_h-h/2\le\OPT_s(A)\le U_h,
\]
so $M=\lceil1/(2\varepsilon)\rceil$ suffices for additive error at most
$\varepsilon T$. These lower inequalities are non-strict. With no switches,
the grid and continuous instance optima agree exactly. To keep a single
uniform interface, the implementation uses the general
certificate~\eqref{eq:coarsening-certificate} throughout.

The lower endpoint may be replaced by zero when $U_h-h<0$, but then the
inequality at that endpoint is non-strict. If $U_h-h=0$, strictness is
retained. This is an additive scheme measured against $T$, not a relative
approximation scheme: an instance optimum can be zero. If the coarse grid
is a subset of a specified fine switching grid $\mathcal T_f$, the same
certificate also brackets the fine optimum, since
\[
 U_h-h<\OPT_s(A)\le\OPT_s^{\mathcal T_f}(A)\le U_h.
\]
For nonnested grids, the coarse schedule remains feasible for continuous
switching, but need not be feasible for the given fine switching grid.
The lower certificate uses switch-preserving transfer and thus does not
extend automatically to dwell-constrained or transition-constrained optima.

\begin{example}[Dwell can prevent convergence under refinement]\label{ex:dwell-grid-gap}
Take $T=1$, two modes, one switch, and minimum dwell $1/2$ for both modes.
Let the input be pure mode one before $1/2$ and pure mode two afterward.
The identical continuous schedule is feasible and has error zero. On a
uniform grid with an odd number $M$ of cells, any feasible two-run schedule
would require both runs to have length exactly $1/2$, which is not a grid
boundary. Only constants are feasible, and their optimum error is $1/2$.
Thus even along arbitrarily fine odd grids the dwell-constrained gap remains
$1/2$. Under these constraints the transfer conclusion itself fails, not
merely one construction.
\end{example}

\begin{corollary}[Input perturbation]\label{cor:coarse-perturbation}
Suppose the data instead describe $\widetilde A\in\Acal_n(T)$ with a known
bound $\norm{A-\widetilde A}_\infty\le\delta$. Let $U_h$ and $V_h$ be
the exact coarse optimum and schedule for $\widetilde A$. Then
\begin{equation}\label{eq:perturbed-certificate}
 U_h-h-\delta<\OPT_s(A)\le D(A,V_h)\le U_h+\delta.
\end{equation}
Consequently the returned schedule has suboptimality strictly below
$h+2\delta$. If the pointwise rate error is bounded by $\rho$, one may
use $\delta=\rho T$.
\end{corollary}
\begin{proof}
Proposition~\ref{prop:compactness} gives
$\OPT_s(A)\ge\OPT_s(\widetilde A)-\delta$; combine it with the strict
lower certificate for $\widetilde A$. The schedule upper bound follows
from the triangle inequality. Integrating a pointwise componentwise rate
bound gives the stated cumulative bound.
\end{proof}
Thus exact optimization of a quantized input is distinct from exact
optimization of the original data. As above, clipping a negative lower
endpoint to zero makes that endpoint non-strict. These certificates do not
require knowing the continuous optimum; they return a feasible schedule
together with explicit lower and upper bounds.
